% GENERATED FILE. Edit canonical lessons, then run npm run generate:books.
% canonical-source-hash: fnv1a64:68eed98a3d0941b6
% canonical-lessons: ZH-C453-beibao, ZH-C453-xinglixiang, ZH-C453-sai, ZH-C453-lin, ZH-C453-biaoqian

\chapter{Bags and Cases}
\label{ch:zh-bags-and-cases}
\hlchaptermodality{453}

\begin{chapteropening}
\textbf{By the end of this chapter:} \emph{I can talk about a backpack, a suitcase, stuffing things in, carrying a case and a luggage tag.}

Complete the last lesson of chapter 453: I can talk about a backpack, a suitcase, stuffing things in, carrying a case and a luggage tag.
\end{chapteropening}

\section[bēibāo]{bēibāo — backpack; rucksack}

\label{lesson:ZH-C453-beibao}



\begin{quote}
Before the new one: say the Mandarin for how dark the sky is getting, then the Mandarin for a paper napkin.
\end{quote}

\subsection*{You'll want to know: bēibāo}
\emph{bēibāo} — \textquotedblleft{}backpack; rucksack\textquotedblright{}.

\begin{quote}
\emph{wǒ de bēibāo bú jiàn le} — My backpack has gone missing.
\end{quote}

Its characters come later; for now it is a word for the ear.

\begin{cousinweb}
\emph{bēibāo} is built from \emph{bēi} 'to carry on the back' and \zh{包} 'a bag', both inherited from Middle Chinese, \emph{pwojH} and \emph{paew} (Guangyun, in Baxter's transcription), which glosses the first as the back of the body and the second 'to wrap'; the level-toned bēi 'to carry on the back' is irregular beside it.
\end{cousinweb}

\begin{grammarlens}[title={What you've built}]
One more word for getting around.
\end{grammarlens}

\subsection*{Your turn}
\begin{itemize}
\raggedright
  \item \emph{Say it:} \emph{bēibāo}
  \item \emph{Say it:} \emph{bēibāo}, once more
  \item \emph{Say it:} \emph{wǒ de bēibāo bú jiàn le} — My backpack has gone missing.
  \item \emph{Recall:} say the Mandarin for how dark the sky is getting, then the Mandarin for a paper napkin, then say \emph{bēibāo} again
\end{itemize}

\begin{tcolorbox}[breakable,colback=teal!4,colframe=teal!35!black,title={Before you move on}]
What does bēibāo mean? (\textquotedblleft{}Backpack; rucksack\textquotedblright{}.) Say it once more.
\end{tcolorbox}
\section[xínglixiāng]{xínglixiāng — luggage case; a travel case}

\label{lesson:ZH-C453-xinglixiang}



\begin{quote}
Before the new one: say the Mandarin for a paper napkin, then the Mandarin for a backpack.
\end{quote}

\subsection*{You'll want to know: xínglixiāng}
\emph{xínglixiāng} — \textquotedblleft{}luggage case; a travel case\textquotedblright{}.

\begin{quote}
\emph{wǒ de xínglixiāng fàng zài fángjiān lǐ} — My suitcase is in the room.
\end{quote}

Its characters come later; for now it is a word for the ear.

\begin{cousinweb}
\emph{xínglixiāng} is built from \emph{xíngli} 'luggage' and \emph{xiāng} 'a box', all inherited from Middle Chinese, \emph{haeng}, \emph{liX} and \emph{sjang} (Guangyun, in Baxter's transcription), which glosses \emph{lǐ} as a fruit tree and gives \emph{xíngli} itself under it, and glosses \emph{xiāng} as a box or a basket.
\end{cousinweb}

\begin{grammarlens}[title={What you've built}]
One more word for getting around.
\end{grammarlens}

\subsection*{Your turn}
\begin{itemize}
\raggedright
  \item \emph{Say it:} \emph{xínglixiāng}
  \item \emph{Say it:} \emph{xínglixiāng}, once more
  \item \emph{Say it:} \emph{wǒ de xínglixiāng fàng zài fángjiān lǐ} — My suitcase is in the room.
  \item \emph{Recall:} say the Mandarin for a paper napkin, then the Mandarin for a backpack, then say \emph{xínglixiāng} again
\end{itemize}

\begin{tcolorbox}[breakable,colback=teal!4,colframe=teal!35!black,title={Before you move on}]
What does xínglixiāng mean? (\textquotedblleft{}Luggage case; a travel case\textquotedblright{}.) Say it once more.
\end{tcolorbox}
\section[sāi]{sāi — to stuff in; to squeeze in}

\label{lesson:ZH-C453-sai}



\begin{quote}
Before the new one: say the Mandarin for a backpack, then the Mandarin for a travel case.
\end{quote}

\subsection*{You'll want to know: sāi}
\emph{sāi} — \textquotedblleft{}to stuff in; to squeeze in\textquotedblright{}.

\begin{quote}
\emph{yīfu dōu sāi zài xiāngzi lǐ le} — The clothes are all stuffed into the case.
\end{quote}

Its characters come later; for now it is a word for the ear.

\begin{cousinweb}
\emph{sāi} is inherited from Middle Chinese \emph{sok} (Guangyun, in Baxter's transcription), glossed 'full', 'stopped up' and 'to separate'; today's level-toned sāi is irregular beside it. The dictionary lists a second reading, \emph{sojH}, glossed as a frontier pass.
\end{cousinweb}

\begin{grammarlens}[title={What you've built}]
One more word for getting around.
\end{grammarlens}

\subsection*{Your turn}
\begin{itemize}
\raggedright
  \item \emph{Say it:} \emph{sāi}
  \item \emph{Say it:} \emph{sāi}, once more
  \item \emph{Say it:} \emph{yīfu dōu sāi zài xiāngzi lǐ le} — The clothes are all stuffed into the case.
  \item \emph{Recall:} say the Mandarin for a backpack, then the Mandarin for a travel case, then say \emph{sāi} again
\end{itemize}

\begin{tcolorbox}[breakable,colback=teal!4,colframe=teal!35!black,title={Before you move on}]
What does sāi mean? (\textquotedblleft{}To stuff in; to squeeze in\textquotedblright{}.) Say it once more.
\end{tcolorbox}
\section[līn]{līn — to carry in the hand (by a handle)}

\label{lesson:ZH-C453-lin}



\begin{quote}
Before the new one: say the Mandarin for a travel case, then the Mandarin for to stuff into a bag.
\end{quote}

\subsection*{You'll want to know: līn}
\emph{līn} — \textquotedblleft{}to carry in the hand (by a handle)\textquotedblright{}.

\begin{quote}
\emph{nǐ néng bāng wǒ līn zhège bāo ma} — Could you carry this bag for me?
\end{quote}

Its characters come later; for now it is a word for the ear.

\begin{cousinweb}
\emph{līn} is inherited from Middle Chinese \emph{leng} (Guangyun, in Baxter's transcription), glossed as holding something so it hangs from the hand; today's līn is irregular beside it.
\end{cousinweb}

\begin{grammarlens}[title={What you've built}]
One more word for getting around.
\end{grammarlens}

\subsection*{Your turn}
\begin{itemize}
\raggedright
  \item \emph{Say it:} \emph{līn}
  \item \emph{Say it:} \emph{līn}, once more
  \item \emph{Say it:} \emph{nǐ néng bāng wǒ līn zhège bāo ma} — Could you carry this bag for me?
  \item \emph{Recall:} say the Mandarin for a travel case, then the Mandarin for to stuff into a bag, then say \emph{līn} again
\end{itemize}

\begin{tcolorbox}[breakable,colback=teal!4,colframe=teal!35!black,title={Before you move on}]
What does līn mean? (\textquotedblleft{}To carry in the hand (by a handle)\textquotedblright{}.) Say it once more.
\end{tcolorbox}
\section[biāoqiān]{biāoqiān — label; luggage tag}

\label{lesson:ZH-C453-biaoqian}



\begin{quote}
Before the new one: say the Mandarin for to stuff into a bag, then the Mandarin for to carry in the hand.
\end{quote}

\subsection*{You'll want to know: biāoqiān}
\emph{biāoqiān} — \textquotedblleft{}label; luggage tag\textquotedblright{}.

\begin{quote}
\emph{měi ge xínglixiāng dōu yǒu yí ge biāoqiān} — Every suitcase has a tag.
\end{quote}

Its characters come later; for now it is a word for the ear.

\begin{cousinweb}
\emph{biāoqiān} is built from \emph{biāo} 'a mark' and \emph{qiān} 'a tally', both inherited from Middle Chinese, \emph{pjiew} and \emph{tshjem} (Guangyun, in Baxter's transcription), which glosses the first 'to raise' and the second, after an older dictionary, 'to check'.
\end{cousinweb}

\begin{grammarlens}[title={What you've built}]
One more word for getting around.
\end{grammarlens}

\subsection*{Your turn}
\begin{itemize}
\raggedright
  \item \emph{Say it:} \emph{biāoqiān}
  \item \emph{Say it:} \emph{biāoqiān}, once more
  \item \emph{Say it:} \emph{měi ge xínglixiāng dōu yǒu yí ge biāoqiān} — Every suitcase has a tag.
  \item \emph{Recall:} say the Mandarin for to stuff into a bag, then the Mandarin for to carry in the hand, then say \emph{biāoqiān} again
\end{itemize}

\begin{tcolorbox}[breakable,colback=teal!4,colframe=teal!35!black,title={Before you move on}]
What does biāoqiān mean? (\textquotedblleft{}Label; luggage tag\textquotedblright{}.) Say it once more.
\end{tcolorbox}
